\dish{Shōyu ramen}
\altdish{醤油ラーメン}
\altdish{Ramen, shōyu (chicken stock)}
\pour{4}
\prep{day before: stock + pork (mostly unattended, 3~h) ; serving
  day: finish + assemble (45~min)}
\source{openai-chatgpt}


\stage{Chicken stock (day before)}

\begin{ingredients}
  \ingr{1.5--2}{kg}{chicken carcasses/backs/necks, plus a few wings}
  \ingr{2.5--3}{L}{cold water}
  \ingr{1}{}{onion, halved}
  \ingr{4--5}{cm}{ginger, sliced}
  \ingr{4}{cloves}{garlic, crushed}
  \ingrS{}{}{green parts of 3--4 spring onions}
\end{ingredients}

\begin{recipe}
  \begin{enumerate}

  \item Bring water and chicken slowly toward a simmer, skimming
    foam.

  \item Add onion, ginger, garlic and spring-onion greens. Maintain a
    very gentle simmer 3~h -- no rolling boil.

  \item Strain. Don't salt. Aim for 1.7--2~L, reduce if you have
    more. Cool and refrigerate.

  \end{enumerate}
\end{recipe}


\stage{Chāshū -- braised pork belly (day before)}

\begin{ingredients}
  \ingr{600--800}{g}{boneless pork belly, skin removed}
  \ingr{120}{ml}{soy sauce}
  \ingr{120}{ml}{sake}
  \ingr{80}{ml}{mirin}
  \ingr{25--30}{g}{sugar}
  \ingr{300--400}{ml}{water}
  \ingr{4}{cm}{ginger, sliced}
  \ingr{3}{cloves}{garlic, crushed}
  \ingrS{2}{}{spring onions, roughly chopped}
\end{ingredients}

\begin{recipe}
  \begin{enumerate}

  \item Brown the pork belly fatty side first, no added oil.

  \item Add remaining ingredients; liquid should come halfway to
    two-thirds up the meat, not necessarily submerge it.

  \item Simmer partially covered 1\textonehalf--2~h, turning every
    20--30~min, until tender but still sliceable.

  \item Cool in the braising liquid, then refrigerate both together
    (cold pork slices far more cleanly).

  \end{enumerate}
\end{recipe}


\stage{Tare (shōyu tare, day before)}

\begin{ingredients}
  \ingr{100}{ml}{Japanese soy sauce}
  \ingr{30}{ml}{mirin}
  \ingr{20}{ml}{sake}
  \ingr{1}{tsp}{sugar}
  \ingr{1}{tsp}{rice vinegar (optional)}
\end{ingredients}

\begin{recipe}
  \begin{enumerate}

  \item Combine, bring briefly to a simmer, remove from heat.
    Reserve.

  \end{enumerate}
\end{recipe}


\stage{Eggs (day before)}

\begin{ingredients}
  \ingr{4}{}{eggs}
  \ingr{60}{ml}{soy sauce}
  \ingr{60}{ml}{water}
  \ingr{30}{ml}{mirin}
\end{ingredients}

\begin{recipe}
  \begin{enumerate}

  \item Cook 6\fracH--7~min, transfer immediately to cold water,
    peel.

  \item Marinate in the soy/water/mirin mix at least 1--2~h, ideally
    overnight.

  \end{enumerate}
\end{recipe}


\stage{Dashi and finishing the broth (serving day)}

\begin{ingredients}
  \ingr{10}{g}{kombu}
  \ingr{15--20}{g}{katsuobushi}
\end{ingredients}

\begin{recipe}
  \begin{enumerate}

  \item Lift the solidified fat off the cold stock; reserve it for
    the aromatic oil.

  \item Put the kombu into the cold stock, leave 30--60~min. Heat
    gently, removing the kombu before it boils (around
    60--70~\textdegree{}C).

  \item Once very hot, turn off heat, add katsuobushi, steep 5~min.
    Strain.

  \end{enumerate}
\end{recipe}


\stage{Aromatic oil (serving day)}

\begin{ingredients}
  \ingr{2}{tbsp}{reserved chicken fat}
  \ingr{1}{}{spring onion, finely sliced}
  \ingrS{}{}{a little ginger or garlic, sliced}
\end{ingredients}

\begin{recipe}
  \begin{enumerate}

  \item Warm the fat gently with the spring onion and ginger/garlic,
    3--5~min, until fragrant. Strain.

  \end{enumerate}
\end{recipe}


\stage{Finishing the pork (serving day)}

\begin{recipe}
  \begin{enumerate}

  \item Lift the solidified fat off the braising liquid. Strain and
    taste it -- it can be used to adjust the tare.

  \item Slice the cold pork (4--5~mm). Warm the slices briefly in
    some braising liquid just before serving, or place them directly
    into the hot bowl.

  \end{enumerate}
\end{recipe}


\stage{Assembly}

\begin{ingredients}
  \ingr{500--600}{g}{\hyperref[ラーメン]{fresh ramen noodles}}
  \ingr{}{}{nori, quartered}
  \ingr{}{}{spring onion, finely sliced}
  \ingr{}{}{menma, mushrooms, pak choi, spinach or bean sprouts
    (optional)}
\end{ingredients}

\begin{recipe}
  \begin{enumerate}

  \item Before cooking the noodles, make a test bowl: 10~ml tare +
    ~120~ml broth. Taste and adjust the tare-to-broth ratio
    (25--35~ml per 350~ml broth, depending on stock and soy sauce).

  \item Heat the bowls. Bring the broth almost to a boil.

  \item Cook the noodles separately, drain very thoroughly, put them
    directly into the bowls.

  \item In each bowl: 25--30~ml tare + 1--2~tsp aromatic oil + ~350~ml
    hot broth, then the noodles.

  \item Top with sliced pork, half an egg, spring onion, nori and
    chosen vegetables.

  \end{enumerate}
\end{recipe}


\stage{Consolidated Shopping List}

\begin{itemize}
  \item 1.5--2~kg chicken carcasses/backs/necks/wings
  \item 600--800~g boneless pork belly
  \item 500--600~g fresh ramen noodles (or \hyperref[ラーメン]{make from scratch})
  \item 4--6~eggs
  \item 1~onion
  \item 2~bunches spring onions
  \item 1~head garlic
  \item ~100~g fresh ginger
  \item kombu
  \item katsuobushi
  \item Japanese soy sauce (~500~ml total)
  \item mirin
  \item sake
  \item sugar
  \item nori
  \item menma (optional)
  \item pak choi/spinach/mushrooms/bean sprouts (optional)
\end{itemize}

%\accord{}
